\chapter{Технологическая часть}

В данном разделе описаны средства разработки и основные особенности реализации компилятора.

\section{Средства разработки}
Компилятор реализован на языке \texttt{Python} версии 3.14.
Управление зависимостями и запуск команд проекта выполняются при помощи инструмента управления проектами \texttt{uv}, что обеспечивает воспроизводимый запуск инструментов из окружения проекта~\cite{python,uv}.

Для синтаксического анализа используется \texttt{ANTLR4} и пакет \texttt{antlr4-python3-runtime}.
Для сборки сгенерированного CIL применяется вспомогательный проект .NET, расположенный в каталоге \texttt{build}.
Проверки проекта выполняются средствами \texttt{pytest}.

\section{Фронтенд компилятора}
Грамматика языка описана в формате \texttt{ANTLR4}.
На ее основе сгенерированы лексер и парсер..
После построения конкретного синтаксического дерева класс построителя AST проходит по узлам ANTLR и создает собственные узлы программы.

Узлы AST реализованы как неизменяемые dataclass-структуры.
Каждый узел содержит позицию в исходном тексте, что позволяет формировать сообщения об ошибках с указанием строки и столбца.
Отдельные группы узлов описывают типы, выражения, операторы, определения типов и определения программ.

\section{Семантический анализ}
Семантический анализатор хранит словари пользовательских типов и программ.
Для каждой программы создается объект области видимости, содержащий известные переменные.
Ошибки накапливаются в списке и выдаются пользователю единым исключением компиляции.

В реализации выделены разрешенные типы: бит, битовая последовательность, массив и запись.
Пользовательский тип при анализе заменяется соответствующим resolved-представлением.
Это позволяет сравнивать типы структурно и не зависеть от имени, которым тип был объявлен в исходной программе.

Особое внимание уделяется компонентному доступу.
При обращении к записи статический компонент выбирает поле по индексу.
При обращении к массиву компонент выбирает тип элемента.
При обращении к битовой последовательности компонент выбирает отдельный бит.
Динамический компонент дополнительно анализируется как выражение.

\section{Генерация CIL}
Генератор CIL формирует текстовую сборку.
Сначала выводится заголовок сборки и определения структур для всех обнаруженных record-типов.
Затем для каждой программы формируется статический метод.
В конце создается метод \texttt{Main}, который разбирает аргументы командной строки как целые числа, вызывает первую программу и выводит результат.

\includelisting
{cil_fragment.il}
{Фрагмент сгенерированного CIL для программы вычисления чисел Фибоначчи}

Для генерации переходов используются уникальные метки.
Такой подход применяется в guard-операторах, \texttt{W}-циклах и \texttt{FOR}-циклах.
Арифметические операции транслируются в инструкции \texttt{add}, \texttt{sub}, \texttt{mul}, \texttt{div}; операция возведения в степень реализуется вызовом \texttt{System.Math::Pow}.

Каждая программа Plankalkul отображается на приватный статический метод CIL.
Если у программы есть имя, оно используется как имя метода; иначе применяется идентификатор вида \texttt{P1}.
Параметры \texttt{V} становятся аргументами метода и загружаются инструкциями \texttt{ldarg}, а промежуточные значения \texttt{Z}, константы \texttt{C} и результаты \texttt{R} помещаются в секцию \texttt{.locals init}.
Для локальных массивов после объявления генерируется явная инициализация: размер помещается на стек инструкцией \texttt{ldc.i4}, затем создается массив инструкцией \texttt{newarr}, после чего ссылка сохраняется в локальную переменную.
Переменные класса \texttt{C} технически также являются локальными переменными CIL, но семантический анализатор запрещает их повторное присваивание.

Отображение типов выполняется после разрешения пользовательских типов.
Бит \texttt{0} соответствует \texttt{bool}, битовые последовательности длиной до 32 разрядов -- \texttt{int32}, более длинные последовательности -- \texttt{int64}.
Массивы языка отображаются на массивы CLR, а записи -- на пользовательские value type структуры.
Для каждой уникальной структурной записи генератор создает тип \texttt{Rec\_N} с полями \texttt{f\_0}, \texttt{f\_1} и далее по числу компонентов.
Именованные типы после семантического анализа сравниваются структурно, поэтому две одинаковые записи могут использовать одно и то же представление в CIL независимо от исходного имени типа.

Результаты программ обрабатываются отдельно от обычных локальных переменных.
Если программа возвращает одно значение, в конце метода загружается соответствующая локальная переменная \texttt{R\_<n>} и выполняется \texttt{ret}.
Если результатов несколько, генератор создает синтетическую record-структуру, заводит локальную переменную \texttt{\_ret}, записывает значения \texttt{R0}, \texttt{R1} и последующих результатов в поля этой структуры инструкциями \texttt{stfld}, затем возвращает структуру целиком.
Метод \texttt{Main} является точкой входа сборки, разбирает аргументы командной строки как \texttt{int32}, вызывает программу с минимальным номером и печатает первый результат через \texttt{System.Console::WriteLine}.

При генерации выражений используется стековая модель CIL.
Целочисленные литералы загружаются инструкциями \texttt{ldc.i4}, булевы значения -- как \texttt{ldc.i4.1} или \texttt{ldc.i4.0}, а битовые последовательности предварительно сворачиваются в целое значение.
Арифметические операции напрямую отображаются на \texttt{add}, \texttt{sub}, \texttt{mul} и \texttt{div}.
Сравнения генерируются через \texttt{ceq}, \texttt{clt}, \texttt{cgt}; операции \texttt{!=}, \texttt{<=} и \texttt{>=} выражаются через строгое сравнение и инверсию результата.
Булевы операции используют \texttt{and} и \texttt{or}, унарное отрицание \texttt{\textasciitilde} реализуется сравнением операнда с нулем, а возведение в степень требует преобразования операндов к \texttt{float64}, вызова \texttt{System.Math::Pow} и обратного преобразования результата к \texttt{int32}.

Компонентный доступ требует различать массивы и записи.
Для массива генератор вычисляет индекс и использует инструкции семейства \texttt{ldelem}/\texttt{stelem}; для массива структур при последующем обращении к полю применяется адресная форма \texttt{ldelema}.
Для записи статический компонент выбирает поле \texttt{f\_<n>}, загрузка выполняется через \texttt{ldfld}, а запись -- через \texttt{stfld}.
При присваивании в компонент генератор сначала получает адрес контейнера, затем вычисляет правую часть и только после этого записывает значение в выбранный элемент или поле.

Управляющие конструкции реализуются через метки и условные переходы.
Guard вычисляет условие и выполняет \texttt{brfalse} на метку конца оператора.
Для цикла \texttt{W} с одним guard используется обычная петля с проверкой условия в начале; для общего случая генератор ведет локальный флаг срабатывания, чтобы повторять цикл только тогда, когда на текущей итерации был выполнен хотя бы один guard.
Цикл \texttt{FOR} инициализирует переменную цикла начальным значением, перед каждой итерацией сравнивает ее с конечным значением, выполняет тело, увеличивает переменную на единицу и переходит к следующей проверке.

\section{Интерфейс командной строки}
Интерфейс командной строки реализует три команды:
\begin{itemize}
	\item \texttt{parse} -- разобрать файл и вывести AST;
	\item \texttt{analyze} -- выполнить синтаксический и семантический анализ;
	\item \texttt{compile} -- сгенерировать CIL и при необходимости собрать программу через .NET.
\end{itemize}

Запуск команд выполняется через \texttt{uv run}, например:
\begin{verbatim}
uv run python -m plankalkulsharp analyze
  tests\programs\test_fibonacci.pk
\end{verbatim}

\section*{Вывод по технологической части}
Была реализована система компиляции на языке \texttt{Python}, использующая ANTLR4 для разбора исходного текста, dataclass-структуры для AST, отдельный семантический анализатор и генератор текстового CIL.
Для воспроизводимого запуска команд проекта используется \texttt{uv}.
